\section{Background and Significance}
\label{sec:background}

Diabetes mellitus affects approximately 537 million adults worldwide, a figure
projected to reach 783 million by 2045.\citep{idf-atlas-10th-edition-2021} In
Israel, the national registry documents a steadily growing patient population
and a corresponding burden on the health system,\citep{israel-national-center-disease-control-diabetes-registry-report-2012-2022,
summary-cost-diabetes-israel-2018} concentrated in hospital care:
approximately 40\% of patients admitted to Israeli general hospitals have
diabetes.\citep{hochberg2023inpatient}

In hospital the clinically relevant object is not the diagnosis but the
glycemic state during the admission. Hyperglycemia arising in hospital is
associated with infection, delayed healing, metabolic decompensation and
increased mortality, and occurs in patients with and without documented
diabetes. Its counterpart, hypoglycemia, is acute, harmful, and largely
iatrogenic---a consequence of the insulin used to correct the first problem.
Inpatient glycemic management is therefore a balancing act shaped less by any
single decision than by how monitoring, medication management and insulin
dosing are carried out repeatedly across a stay. That makes the treatment
process itself, not only the patient's state, worth measuring.

In September 2023 the Israel Endocrine Society, the Israel Diabetes Association
and the National Diabetes Council issued a national position paper on inpatient
diabetes care,\citep{hochberg2023inpatient} specifying how often glucose
should be measured, how oral antidiabetic agents should be handled on
admission, when basal and correction insulin should be given, and how to
respond to glycemic episodes. Recommendations of this kind are checkable in
principle because they are procedural, temporal and conditional: an action is
expected, within some window or at some frequency, when a triggering situation
holds. Establishing whether they were followed is harder, having conventionally
required chart review, which does not scale past modest single-institution
samples---too small to separate one ward's practice from the effect of a
guideline, and far too small to observe how practice shifts when a guideline is
introduced nationally.

Informatics supplies the machinery to do this automatically. Knowledge-based
temporal abstraction and the Temporal Mediator \citep{shahar1997framework} turn
time-stamped clinical data into named interval concepts; Asbru
\citep{shahar1997asbru, shahar1998asgaard} expresses a guideline's procedural
content as a hierarchy of execution plans; and Picard
\citep{shalom2016architecture} combines the two to generate real-time
recommendations and adherence reports. Deployed in MobiGuide,\citep{peleg2022mobiguide} it raised clinicians' guideline compliance from
41\% to 93\%---direct evidence that adherence is malleable, and therefore worth
measuring. Retrospective assessment required a third body of knowledge:
\emph{quality-assessment} knowledge \citep{ben2024supporting} re-expresses that
procedural content declaratively, as the temporal pattern correct execution
would have left in the record, with a membership function grading partial
compliance rather than a binary verdict. Such systems have been validated
against simulated scenarios \citep{shalom2015multiple} and applied
institutionally to long-term management at a geriatric hospital.\citep{shalom2024implementation}

This study adopts that formulation of adherence but not the systems that
introduced it, and departs from them in three respects. In
\emph{representation}, procedural and declarative knowledge are encoded in a
single schema evaluated by one engine, and the resulting knowledge base is
itself a transferable contribution. In \emph{scale}, adherence is measured
across six hospitals spanning a national implementation rather than within one
institution. In \emph{purpose}, prior work evaluates care \emph{against} the
guideline, which serves as the reference standard; here each recommendation is
treated as a hypothesis to be tested against what subsequently happened to the
patient. Adherence thus becomes an explanatory variable: the question is not
whether the guideline was applied correctly, but which of its recommendations
demonstrably matter.

\section{Objective}
\label{sec:objective}

We set out to determine, across a national hospital network spanning the
introduction of an inpatient hyperglycemia guideline, how completely each of
its individual recommendations was implemented and whether that changed after
the guideline was issued; whether adherence is associated with subsequent
complications once measured patient characteristics are accounted for; and
which individual recommendations carry that association.
